\section{Introduction}
\label{sec:introduction}

Scholarly metadata is spread across many bibliographic providers. Open indexes such as OpenAlex expose works, authors and institutions at scale~\cite{priem2022openalex}, but the providers differ in what they cover, and coverage comparisons show that the differences depend on the task and on the sources compared~\cite{culbert2025openalexcoverage}. A record for the same publication may therefore be incomplete, formatted differently, or duplicated across sources, and its author list may be abbreviated, reordered or attached to different identifiers. An application that aggregates such records has to decide two separate questions.

The first is whether two records describe the same \emph{work}. This is a record-linkage and deduplication problem with a long statistical and practical history~\cite{fellegi1969recordlinkage,elmagarmid2007duplicate,forbes2024deduplicator}. The second is whether a name on a record, or a name typed by a user, refers to a particular \emph{person}. This is author name disambiguation, a mature field with several task formulations and a large literature~\cite{ferreira2012brief,rodrigues2024andreview,debonis2023graph}. The two questions use different evidence and carry different costs of error. A shared DOI is strong evidence that two records are the same work, whereas two researchers with the same or similar names may be different people, and one researcher may appear under several name forms. Treating a work-level match as if it settled a person-level question, or the reverse, hides the uncertainty that a user most needs to see.

This paper concerns the second question, in the setting of a scholarly metadata aggregator. We use the following terms consistently. \emph{Candidate retrieval} obtains a set of possible author profiles for a query. \emph{Ranking} orders that set. An \emph{identity decision} is the choice among three outcomes: automatic assignment of one candidate, manual review, or an unresolved result. These steps are not interchangeable. A retriever that never returns the correct profile cannot be rescued by any ranker or decision rule, and a ranker that places the correct profile first does not by itself justify assigning it automatically.

The literature treats the surrounding pieces separately. Benchmark studies evaluate clustering or assignment of publications to authors, typically with F1-style measures, and report that models trained on one dataset may generalize poorly to others~\cite{subramanian2021s2and,chen2023whoiswho}. General work on selective prediction supplies the vocabulary of coverage and risk for systems that may decline to answer~\cite{chow1970reject,elyaniv2010foundations,geifman2019selectivenet}. Deployed systems combine automatic disambiguation with human annotators~\cite{zhang2018aminerhitl}. Within the sources we reviewed, we did not identify a study that reports the retrieval stage and the selective assignment stage together, on the same frozen candidate pool, with manual-review, unresolved and provider-failure outcomes kept in the denominators. This is a statement about the limited search described in Section~\ref{sec:related}, not a claim that no such study exists.

We study this problem through \hammer{} (Hybrid Academic Metadata Aggregator and Evidence-Guided Author Entity Resolution System), an implemented application. \hammer{} aggregates metadata from several providers, normalizes and deduplicates publication records, and stores provenance for each source record. Its author-resolution component retrieves candidate author profiles from OpenAlex and scores them with five evidence signals: name similarity, institutional evidence, publication overlap, coauthor evidence and temporal compatibility. The signal weights are engineering defaults and are renormalized over the signals that are available; the resulting score ranks candidates and feeds a set of conservative decision gates, but it is a heuristic score and not a calibrated probability. When evidence is missing or a provider call fails, the system reports that state explicitly and does not treat it as evidence against a candidate.

\hammer{} does not propose a new disambiguation algorithm, and we make no claim of algorithmic novelty. Each ingredient it uses has prior art, and combining existing components does not by itself constitute a contribution to the method. The contribution we aim to establish is narrower: a provenance-aware workflow that keeps work-level matching separate from person-level decisions, together with an evaluation protocol that measures how its decision gates trade automatic precision against coverage. Our primary research question is:

\begin{quote}
\textbf{RQ2.} How effectively does \hammer{} retrieve author candidates, and how do its evidence-guided decision gates balance automatic assignment precision against coverage?
\end{quote}

We separate it into \textbf{RQ2a}, which asks how effectively the candidate retrieval process returns the expected author profile (Recall@$k$, mean reciprocal rank, candidate pool size, empty-result rate and candidate-pool mapping rate), and \textbf{RQ2b}, which asks how the decision gates affect the trade-off between the correctness of automatic assignments and automatic coverage (automatic precision, automatic coverage, manual-review and unresolved rates, selective risk and false automatic assignments). Precision and coverage are always reported together: a system that assigns few cases automatically cannot be called accurate without that coverage being visible, and a system that abstains on every case is not safe merely because it makes no wrong assignment.

The contributions of this paper, subject to the evidence presented in the later sections, are:
\begin{enumerate}
  \item An integrated architecture for scholarly metadata aggregation with source-level provenance, in which publication-level deduplication and author-level identity decisions are separate stages (Section~\ref{sec:architecture}).
  \item A staged author-resolution workflow, from candidate retrieval through evidence extraction and scoring to conservative decision gates, with explicit handling of missing and failed evidence and three outcomes: automatic assignment, manual review, and unresolved (Section~\ref{sec:methodology}).
  \item An evaluation protocol that reports retrieval effectiveness and selective resolution on one frozen candidate pool, controls for label and evidence leakage, and keeps uncertain and unresolved cases in the record (Section~\ref{sec:design}).
  \item Reproducible experimental artifacts: byte-level capture of provider responses with recorded hashes, and offline replay of the resolver from those captures (Section~\ref{sec:design}).
  \item Exploratory, benchmark-derived observations on 68 real queries that locate where the gates lose coverage (competition among similarly named candidates, not low scores) and that show the heuristic ordering of a pool doing no better than the provider's own order (Section~\ref{sec:results}).
\end{enumerate}

This manuscript reports a complete system and a fully specified evaluation, but not a verified accuracy. The independent identity-label audit on which the primary claims depend has not been performed, and no precision, recall or coverage figure is claimed as a property of the system. We report instead exploratory observations from 68 real queries (Section~\ref{sec:results}), labelled as benchmark-derived, small and clustered, and we state what they do and do not show.

The remainder of the paper is organized as follows. Section~\ref{sec:related} reviews related work and positions \hammer{} against it. Section~\ref{sec:architecture} describes the system architecture, Section~\ref{sec:methodology} the author-resolution method, and Section~\ref{sec:design} the experimental design and evaluation protocol. Sections~\ref{sec:results}--\ref{sec:conclusion} cover results, discussion, threats to validity, and conclusions.
